% Requires graphicx. Paths are relative to a manuscript in paper/.
% The middle figure is an adapted control design, not an implemented method.
\begin{figure*}[t]
  \centering
  \includegraphics[width=\textwidth]{figures/loop-residual-comparison-v3/vanilla-loop.pdf}
  \caption{Vanilla Loop, retained as an overall architectural reference.
  The shared Transformer's raw output is passed directly to the next iteration.
  This reference does not use the matched soft-feedback candidate of the other two figures.}
  \label{fig:matched-highway-vanilla}
\end{figure*}

\begin{figure*}[t]
  \centering
  \includegraphics[width=\textwidth]{figures/loop-residual-comparison-v3/highway-loop.pdf}
  \caption{Highway-gated Loop (adapted), a proposed mechanism control.
  A shared linear-plus-sigmoid gate reads the previous recurrent state and
  mixes it with the new candidate. Candidate computation, input reinjection,
  soft feedback and RMS normalization match Invariant Loop. Both the Transformer
  and gate parameters are reused across loops. The gate follows Highway Networks;
  the shared-Transformer adaptation is not an architecture specified by that paper.
  Each displayed token has a local carry; only one gate is expanded.}
  \label{fig:matched-highway-control}
\end{figure*}

\begin{figure*}[t]
  \centering
  \includegraphics[width=\textwidth]{figures/loop-residual-comparison-v3/invariant-loop.pdf}
  \caption{Invariant Loop, following the implementation. The common candidate
  path adds soft feedback and normalizes the result. Five statistics feed an MLP;
  historical-mass normalization yields scalar $\gamma_t$ for each token.
  The two aggregation diagrams differ in the aggregation module as a whole,
  including gate input, dimensionality and normalization. They do not establish
  separate causal contributions of these choices.}
  \label{fig:matched-highway-invariant}
\end{figure*}
